\documentclass[11pt]{article}
\input{preamble}
\title{R02 --- Datatypes and recursion \\ \large Reference solutions (author-only)}
\author{dafny-verus-learn-0}
\date{}
\begin{document}
\maketitle
\section*{Problem B1}\label{ans:B1}
\begin{answer}
\textbf{Task shape:} The card describes 161 HumanEval tasks re-specified in Lean. A system supplies a specification with a faithfulness proof, and an implementation with its correctness proof.

\textbf{Oracle:} Specification certification compares the candidate and hidden reference for all implementations; implementation correctness is judged against the hidden reference, not merely the generated specification.

\textbf{Trust:} The result relies on Lean and its logical foundation plus the faithfulness of the human reference and task translation. Successful checking does not certify the natural-language intent without that external reference judgement.

\textbf{Cheat surface:} A directly executable specification can be copied into the implementation, while a weak property may admit a constant program. The curated non-computable specifications and equivalence obligation target those shortcuts.

\textbf{Metric:} The source reports stage-wise compilation and proof rates plus end-to-end success under pass-at-600-seconds. The largest reported end-to-end count is one task out of 161; a high compilation rate alone is not end-to-end success.

\textbf{Budget:} The reported allowance is 600 seconds; the card also states approximately 25 minutes of specification authoring plus 15 minutes of review. It supplies no common token or model-call budget for all systems.

\textbf{Variance:} The range reports a severe floor effect and no error bars. It does not provide repeated-run uncertainty estimates from which to infer a stable ranking between nearly tied systems.

\textbf{Artifact:} The named artifact is the curated Lean benchmark and its reference specifications. The source range does not supply a release commit, download checksum or licence; those would need separate evidence.
\end{answer}
\section*{Problem B2}\label{ans:B2}
\begin{answer}
The measure is $(r,|s|)$ in lexicographic order. An Alt or Cat call
uses a strict constructor child for the first component, independently of the
slice length. In Star's first recursive call the regex is its strict child.
In its second call the regex is the same Star value, but $0<i\leq |s|$ makes
$|s[i..]|=|s|-i<|s|$. Allowing $i=0$ loses that strict second-component decrease.
Removing the length component loses the comparison for the second Star call
even when a nonempty prefix is consumed. The argument needs both the bound
and the lexicographic order, not the recursive call's spelling alone.
\end{answer}
\section*{Problem B3}\label{ans:B3}
\begin{answer}
For \texttt{r}, the only successful split is $i=0$: the prefix is
empty and uses Eps in Alt; the suffix is \texttt{r} and uses its single Chr.
For \texttt{qr}, the only successful split is $i=1$: prefix \texttt{q} uses
Alt's Chr, and suffix \texttt{r} uses the final Chr. For \texttt{rr}, $i=0$
leaves a two-character suffix, which the final Chr cannot match; $i=1$ gives
prefix \texttt{r}, which is neither the empty string nor \texttt{q}; $i=2$
leaves an empty suffix, also not matched by the final Chr. These are all
indices from zero through two. Each constructor consumes its entire slice;
finding an occurrence of the final character inside a longer slice is insufficient.
\end{answer}

\end{document}
